%\vspace{-1.5ex}
\section{Conclusion}
In this paper, we propose a novel framework \textbf{K-RagRec} to augment the recommendation capability of LLMs by retrieving reliable and up-to-date knowledge from KGs. Specifically, 
we first introduce a GNN and PLM to perform semantic indexing of KGs, enabling both coarse and fine-grained retrieval for KGs. To further improve retrieval efficiency, we introduce a popularity selective retrieval policy that determines whether an item needs to be retrieved based on its popularity. 
Notably, K-RagRec performs more expressive graph encoding of the retrieved knowledge sub-graphs, facilitating the LLM to effectively leverage the structure information and avoid long context input. Extensive experiments conducted on three real-world datasets demonstrate the effectiveness of our proposed framework.







\section{Limitations}
To the best of our knowledge, this work is a pioneering study in investigating knowledge-graph RAG for LLM-based recommendations.  Therefore, we realise that this work still has the following three main limitations that can be improved in the future research:
\vspace{-0.1in}
 \begin{itemize} [leftmargin=*] 
\setlength{\itemsep}{0pt}
\setlength{\parsep}{0pt}
\setlength{\parskip}{0pt}
\item Firstly, due to the GPU resource constraints, we were only able to evaluate our framework on 7b and 8b models. Therefore, in future work, we plan to extend our method to larger models to fully assess its effectiveness and scalability.
\item Additionally, we only utilize Freebase as the external KG as it is most commonly used for recommendation tasks. 
Thus we also aim to adopt other KGs, such as YAGO, DBpedia, and Wikipedia, to better understand how different knowledge sources impact the performance of the proposed method.
\item Lastly, designing an intelligent selective retrieval policy for LLM-based recommender systems is an important and challenging task.
In this work, we propose to leverage popularity to
determine which items to retrieve to improve retrieval efficiency.
In the future, we will investigate more flexible mechanisms (e.g., reinforcement learning) to dynamically update policies according to changes in user interest. 
\end{itemize}

\section{Acknowledgements}
The research described in this paper has been partly supported by General Research Funds from the Hong Kong Research Grants Council (project no. PolyU 15207322, 15200023, 15206024, and 15224524), internal research funds from The Hong Kong Polytechnic University (project no. P0042693, P0048625, P0051361, P0052406, and P0052986). 

